\begin{table}[!h]
\centering
\caption{\label{tab:gp_summary_combined}Summary of Gram Panchayats by Election Year}
\centering
\fontsize{8}{10}\selectfont
\begin{threeparttable}
\begin{tabular}[t]{lrrrr}
\toprule
 & 2005 & 2010 & 2015 & 2020\\
\midrule
\addlinespace[0.3em]
\multicolumn{5}{l}{\textit{\textbf{Panel A: Rajasthan}}}\\
\hspace{1em}Gram Panchayats & 9,178 & 9,166 & 9,862 & 11,506\\
\hspace{1em}Panchayat Samitis & 237 & 249 & 300 & 352\\
\hspace{1em}Districts & 32 & 33 & 33 & 33\\
\hspace{1em}Women Reserved (\%) & 33.5 & 47.6 & 48.4 & 50.0\\
\addlinespace[0.3em]
\multicolumn{5}{l}{\textit{\textbf{Panel B: Uttar Pradesh}}}\\
\hspace{1em}Gram Panchayats & 51,737 & 41,227 & 58,994 & 49,772\\
\hspace{1em}Panchayat Samitis & 846 & 844 & 823 & 728\\
\hspace{1em}Districts & 70 & 72 & 75 & 67\\
\hspace{1em}Women Reserved (\%) & 50.0 & 50.0 & 50.0 & 50.0\\
\bottomrule
\end{tabular}
\begin{tablenotes}
\item \textit{Note: } 
\item GP counts represent the total number of Gram Panchayats for which we have electoral data. The number of GPs increased between 2010 and 2020 due to delimitation (redistricting). The 2020 column for UP contains 2021 election data.
\end{tablenotes}
\end{threeparttable}
\end{table}
